\documentclass[a4paper]{article}

\usepackage[english]{babel}
%\usepackage[utf8x]{inputenc}
\usepackage{amsmath}
\usepackage{graphicx}
\usepackage[colorinlistoftodos]{todonotes}
\usepackage[all]{xy}
\usepackage{dsfont}
\usepackage{amssymb}
\usepackage{bbold}
\usepackage{amsthm}

\usepackage{tikz}
\usepackage{verbatim}
\usetikzlibrary{matrix}
\usetikzlibrary{cd}

\usepackage[autostyle=true]{csquotes}
\usepackage{pdfpages}
\usepackage{hyperref}
\usepackage{fancyhdr}

\newtheoremstyle{exampstyle}
  {\topsep} % Space above
  {\topsep} % Space below
  {\itshape} % Body font
  {} % Indent amount
  {\bfseries} % Theorem head font
  {.} % Punctuation after theorem head
  {.5em} % Space after theorem head
  {} % Theorem head spec (can be left empty, meaning `normal')
\theoremstyle{exampstyle}
\newtheorem{thm}{Teorema}[chapter]
\newtheorem{prop}[thm]{Proposici\'on}
\newtheorem{lemma}[thm]{Lemma}
\newtheorem{cor}[thm]{Corolario}
\theoremstyle{definition}
\newtheorem{thme}{Expected Theorem}

\newtheorem{rec}[thm]{Recuerdo}
\newtheorem{conj}{Conjecture}
\newtheorem{defin}[thm]{Definici\'on}


\newcommand{\norm}[1]{\left\lVert#1\right\rVert}

\newcommand{\Rn}{\mathbb{R}^n}


\newcommand{\too}[1]{\overset{#1}{\longrightarrow}}
\newcommand{\hookar}[1]{\overset{#1}{\hookrightarrow}}
\newcommand{\colim}[1]{ \underset{#1}{\texttt{colim}} \text{ } }
\newcommand{\fset}[1]{ \{ #1_0 , \dots , #1_n \} }
\newcommand{\ddt}{\frac{d}{dt}}


\newcommand{\N}{\mathbb{N}}
\newcommand{\Z}{\mathbb{Z}}







\fancyhead[L]{\includegraphics[scale=0.1]{UFT.png}}

\begin{document}


\thispagestyle{fancy}
\vspace*{1cm}

\begin{minipage}[r]{12cm}
    \begin{flushright}
\begin{tikzpicture}
\draw[green, very thick] (0,0) rectangle (3,2);
\end{tikzpicture}
    \end{flushright}
\end{minipage}

\begin{center}
\Large{\textbf{Control recuperativo\\ Álgebra Lineal\\ }}
\end{center}
\vspace{1cm}
\begin{minipage}[r]{15cm}
\begin{flushleft}

Fecha: 22 de noviembre 2024 \\%
Duración: 120 minutos  \\
Profesores: Camila Muñoz.\\
\end{flushleft}
\end{minipage}
\vspace{2cm}
\begin{center}
\Large{\textbf{Facultad de Ingeniería}}
\end{center}
\vspace{1cm}
\textbf{En caso de que un alumno sea sorprendido en actitudes deshonestas, ya sea intentando o
dejándose copiar, será calificado con nota mínima 1.0 y perderá toda posibilidad de eliminar
dicha nota. Además, esta conducta quedará archivada como antecedente en su carpeta y
podrá ser causal de eliminación.
Así mismo, está prohibido salir de la sala, usar o mantener sobre la mesa teléfonos celulares
o cualquier tipo de aparato electrónico, a no ser que el profesor lo requiera.}\\



\vspace{5cm}
\makebox[12cm]{\textbf{Nombre Alumno}:\ \hrulefill} 
\newpage
\thispagestyle{empty}
\section*{Preguntas}

\begin{enumerate}
\item \textbf{(10 puntos)} Usando el teorema del rango, identifique si el siguiente sistema 
    \begin{align*}
x + y +2z &= 0 \\
2x +y -  z &= 3\\
3x +2y +  z &= 3
\end{align*}
tiene una, infinitas o ninguna soluci\'on. En caso de existir, encuentre el conjunto soluci\'on.

\item Considere el conjunto de vectores $B = \left\{ ( 1 , 0 , 0) , ( 2 \, 1 , 0 ) , (1 , 1 , 1 ) \right\}$ en $\mathbb R^{3}$. 
\begin{enumerate}
    \item \textbf{(5 puntos)} Verifique que $B$ es base de $\mathbb R^{3}$. 
    \item \textbf{(10 puntos)} Encuentre una base ortonormal de $\mathbb R^{3}$ a partir de $B$ usando el método de Gram-Schmidt.
    \item \textbf{(5 puntos)} Encuentre las coordenadas del vector $  (-1, 0 , 1 ) $ en la base obtenida en la parte (b). 
\end{enumerate}


\item Considere la funci\'on 
      $T:\mathbb{R}^3\to P_2[x]$, definida por $$T(a,b,c)=(2a+b)+(a-b)x+ (a-c)x^2$$
      \begin{enumerate}
          \item \textbf{(5 puntos)} Muestre que $T$ es una transformaci\'on lineal. 
          \item \textbf{(5 puntos)} Encuentre el kernel de $T$, una base y su dimensi\'on \item \textbf{(5 puntos)} Determine la dimensi\'on de la imagen.
          \item \textbf{(5 puntos)} Encuentre la matriz de la transformaci\'on lineal en las bases can\'onicas.
      \end{enumerate}
    \item \textbf{(10 puntos)} Considere la siguiente matriz
\[    A=\begin{pmatrix} 1 &2 & 1 \\ 0 & -1 & -2 \\ 0 & 0 & 3 \end{pmatrix}.\]
Encuentre los valores y vectores propios de $A$.
    

\newpage
        
\section*{Formulario} 
\begin{itemize} 
\item \textbf{Bases:} Un conjunto de vectores $\{ v_1, \dots , v_n \}$ se dice linealmente independiente si $ \lambda_1 v_1+ \dots + \lambda_n v_n= \Vec{0} \Longrightarrow \lambda_1= \dots = \lambda_n= 0$. $W$ se dice generado por $\{w_1, \dots , w_n\}$ si $<\{w_1, \dots , w_n\}>=\{ \sum \lambda_iw_i | \lambda_i \in \mathbb{R} \}=W$ Una base de un espacio vectorial $V$ es un conjunto linealmente independiente que genera $V$. \item \textbf{Procedimiento para obtener una base ortonormal:} Sea $B = \{v_i\}_{i=1}^m $ una base ordenada cualquiera de $W $. \begin{itemize} \item Escribimos $w_1 = v_1$. \item Calculamos $w_2=v_2-\dfrac{<v_2,w_1>}{<w_1,w_1>}w_1$. \item Calculamos $w_3=v_3-\dfrac{<v_3,w_1>}{<w_1,w_1>}v_1- \dfrac{<v_3,w_2>}{<w_2,w_2>}w_2$. \item Inductivamente obtenemos los vectores $w_m=v_m-\sum_{i=1}^{m-1}c_iw_i$ con $c_i=\frac{<v_m,w_i>}{<w_i,w_i>}$. \end{itemize} El conjunto $\{ w_1,w_2, \dots, w_n \}$ es una base ortogonal. Si se normaliza cada vector $w_k$, obtenemos nuevos vectores $u_k=\frac{w_k}{\|w_k\|}$. Luego, $\{ u_1,u_2, \dots, u_n \}$ una base ortonormal para $W$. 
\item \textbf{ Coordenada respecto de una base ortogonal } $B = \{w_i\}_{i=1}^m $ est\'a dada por $a_i=\dfrac{\langle v, v_i \rangle}{\langle v_i, v_i \rangle}$ \item {\bf Proyección ortogonal: } Sea $W$ un subespacio de $\mathbb{R}^n$ con base ortogonal $\{ w_1,w_2, \dots , w_k \}$. Si $v \in \mathbb{R}^n$, entonces la proyección ortogonal de $v$ sobre $W$, denotada por $P_W(v)$, está dada por: $$P_W( v) = \frac{\langle v, w_1 \rangle}{\langle w_1, w_1 \rangle}w_1 + \frac{\langle v, w_2 \rangle}{\langle w_2, w_2 \rangle}w_2 + \cdots + \frac{\langle v, w_k \rangle}{\langle w_k, w_k \rangle}w_k.$$ \item {\bf Transformaciones lineales: } Transformación Lineal $T$ de $ V$ en $W$, como una función $T:V\to W$ tal que \begin{enumerate} \item $T(v+w)=T(v)+T(w)$, \item $T(\alpha v)=\alpha T(v)$, \end{enumerate} para todo $v,w\in V$, para todo $\alpha \in K$. \item $T:V\to W$ transformaci\'on lineal, entonces $$\operatorname{dim}V=\operatorname{dim} \operatorname{ker }T+\operatorname{dim}\operatorname{Im}T$$ \item {\bf Valores y Vectores Propios: } $\lambda \in \mathbb{R}$ es valor propio de $M \in \mathcal{M}_{n \times n}$ si existe $\Vec{v}\in \mathbb{R}^n$ con $M\Vec{v}= \lambda \Vec{v}$. $\Vec{v}$ se llama vector propio asociado a $\lambda$ \begin{itemize} \item Los valores propios son las raices del polinomio característico $p(x)=\texttt{det}(M - x I_n ) $ \item Los vectores propios asociados a un valor propio $ \lambda_0$ forman el subespacio propio $E_{\lambda}=\texttt{Ker}(M-\lambda_0 I_n)$ \item La multiplicidad geométrica de $\lambda$ es $\texttt{dim}(E_\lambda)$ y la multiplicidad algebráica es su multiplicidad como raiz del polinomio característico. \end{itemize} \end{itemize}






\end{enumerate}


\end{document}